% 13-operations.tex: surviving operational reality.
\section{Surviving reality: the stack, failure, load and the standing witness}
\label{sec:operations}

\subsection{The stack}

The production topology is five services: a self-built TLS edge that negotiates a hybrid
post-quantum key exchange with modern clients, the application under gunicorn, a connection pooler,
PostgreSQL with continuous archiving, and Redis for the shared rate limiter. Every service runs as a
non-root user with every capability dropped, third-party images are pinned by digest, and every
production container drops the test framework. Four deployment paths share one secrets discipline,
one backup path and one threat model: a laptop launcher, a single-host compose profile with a
blue-green rolling deploy, a scripted Linux install under systemd exercised on two distributions,
and a Helm reference profile with default-deny network policies and the restricted pod security
standard, booted in CI on a one-node cluster. \sImpl

\subsection{Failure, drilled and measured}

The repository's rule is to exercise rather than read: ten of ten defects found in one sweep were
invisible to reading. Every operational claim below is a drill that runs in CI and commits a number.

\begin{description}
  \item[High availability.] The database runs under Patroni with a leader lease in a three-member
    etcd and role-based routing. Under a live write stream the drill loses the leader (the replica
    promoted after the 20-second lease, in-flight queries failing fast and retried, none lost), cuts
    the leader off from the lease store (it stood down in 7 seconds), performs a planned switchover
    (3.3 seconds, no insert failed), and crashes an etcd member (a 0.3-second stall). The Kubernetes
    profile runs the same members with the cluster's API as the lease store and deletes the leader
    pod. \sImpl
  \item[Read replicas.] Read-only surfaces route to a streaming replica under an explicit
    staleness contract, at most ten seconds behind by default, falling back to the primary beyond
    it, with every routed response naming its source and lag. Correctness-critical reads stay on
    the primary, and the authorization half of a verification verdict is pinned there
    (Section~\ref{sec:authorization}). \sImpl
  \item[Partitioning.] The four append-only event tables are monthly range-partitioned, and a
    drill proves that the append-only trigger holds across a partition, an attach and a detach,
    because the archives' carve-out must not become a rewrite path. A benchmark found that the
    Atlas roll-ups defeated partition pruning under the generic plan and scanned every month of
    history; the fix is pinned by a check. \sImpl
  \item[Disaster recovery.] Backups are encrypted and a restore fails closed without the key.
    A monthly drill kills the primary, restores it from the archive and commits the measured numbers
    against the targets of 300 seconds of data loss and four hours to recover: the last recorded
    rows show 36 to 42 seconds of loss and under five seconds to service, on the drill's ephemeral
    host, not production hardware. \sImpl
  \item[Chaos.] A weekly drill boots the blue-green stack under traffic, kills one colour with zero
    dropped requests, stops both until the outage page reaches a real alert manager and webhook,
    kills Redis and PostgreSQL, partitions the pooler, and commits every recovery time against a
    ceiling. Findings become checks. \sImpl
  \item[Rolling deploys.] A deploy under traffic drops zero requests; edge and database recreation
    are window operations measured on every push against ceilings (0.3 seconds for the edge on a
    single host, 0.6 seconds for a database crash), and the readiness ledger carries the remaining
    edge half openly. \sImpl
\end{description}

\subsection{Load: what was measured and what was projected}

Three kinds of number appear in the repository, and the paper keeps them apart. \emph{Measured}
numbers come from a named machine, an Apple Silicon laptop with eight cores, and carry a version
stamp. \emph{Simulated} numbers come from the national simulation harness, a synthetic country
driven through the real procedures at a stated scale factor. \emph{Projected} numbers are arithmetic
over measured ones and are labelled as such by the repository itself. No production cluster has
been measured.

{\footnotesize
\setlength{\tabcolsep}{4pt}\renewcommand{\arraystretch}{1.12}
\rowcolors{2}{tabstripe}{white}
\begin{xltabular}{\textwidth}{>{\RaggedRight\arraybackslash}X >{\RaggedRight\arraybackslash}p{1.8in} >{\RaggedRight\arraybackslash}p{0.8in} >{\RaggedRight\arraybackslash}p{1.15in}}
\caption{Performance numbers at the versions that measured them. Every row is single-node on the reference laptop unless marked as a projection. The event-ingestion rate and the cryptographic verification rate differ by about 35 times, a gap the repository's own history once hid (Section~\ref{sec:senses}).}\label{tab:perf}\\
\toprule
\rowcolor{tabheader}\thead{Quantity} & \thead{Number} & \thead{Kind} & \thead{Source} \\
\midrule
\endfirsthead
\toprule
\rowcolor{tabheader}\thead{Quantity} & \thead{Number} & \thead{Kind} & \thead{Source} \\
\midrule
\endhead
\midrule\multicolumn{4}{r}{\itshape continued on the next page}\\
\endfoot
\bottomrule
\endlastfoot
ML-DSA-65 sign, one core & about 2,110 per second & measured & dyno, \code{v9.278} \\
ML-DSA-65 verify, single witness, one core & about 7,840 per second & measured & dyno, \code{v9.278} \\
ML-DSA-65 verify, two witnesses, one core & about 740 per second & measured & dyno, \code{v9.278} \\
Membership proof, depth 14 & about 35 ms to prove, 13 ms to verify, 76 KB & measured & dyno and the soundness ledger \\
Issuance through the application & 40 per second sustained, p95 28 ms & measured & baseline, \code{v9.191} \\
Verification event through the application & 80 per second sustained, p95 18.5 ms & measured & baseline, \code{v9.191} \\
Enrolment with real signing, bulk pipeline & about 372 tokens per second & simulated, 1:10,000 & national simulation, \code{v9.257} \\
Verification-event ingestion (audit writes, no signature check) & about 25,970 per second & simulated, 1:10,000 & national simulation \\
Atlas at ten million events: a street block; the whole world; the whole world from a roll-up & 2.6 ms; 2.9 s; 0.04 ms & measured, \code{v9.150} & scaling ledger \\
Single-witness verification across an eight-core node & about 62,800 per second & projection & benchmark arithmetic \\
\end{xltabular}
}

The roadmap's planning targets (ten million persons for a state; 350 million, five thousand
sustained and fifty thousand peak verifications per second nationally) are targets to be validated,
stated in the repository as such and not asserted. A national workload has been simulated at a
scale factor and driven through the real constraints; it has not been run. \sExp

\subsection{Abuse, secrets and supply chain}

Per-agency quotas on issuance, revocation and verification are opt-in rows enforced by a trigger on
every write path, exact under concurrency by an advisory lock, refused loudly on every surface, and
keyed by agency identifiers, never by a person. Velocity alerts compare each agency to its own
trailing week. The per-IP rate limiter shares one Redis bucket across workers. Operators
authenticate with WebAuthn as a second factor after an enrolment deadline, sessions are registered
server-side with per-role caps and origin checks, and a per-role network policy applies. Secrets are
file-mounted, sealed and rotated through one lifecycle; the application refuses to boot in
production on a default key. Every release ships bills of materials with signed provenance, and CVE
scans of the dependency surface and of every self-built image gate the build. \sImpl

\subsection{The standing witness}

Eighteen CI jobs run on every push, and most carry a comment naming the past failure they exist to
prevent. They build and boot the images; boot the five-service stack end to end and assert health
through the TLS edge; round-trip an encrypted backup and an off-site copy; kill and restore the
primary; roll a deploy under traffic; install the Linux profile on two distributions; boot the
Kubernetes profile; page a duress alert through a real alert manager; prove the post-quantum
handshake against a real certificate; sign and verify with real ML-DSA both in the production image
and inside a software PKCS\#11 module; run every protocol drill under real signatures; federate two
instances over HTTP; type-check, test and conform the TypeScript SDK; fail the HA profile's leader,
lease store and etcd member under a live write stream; and gate on CVE scans. The pattern behind
them is the one Section~\ref{sec:senses} describes: every defect class found by running the system
gets a permanent gate. \sImpl

\begin{layercard}{Layer card: operations}
\qa{What it is}{The five-service stack on four substrates; HA, replicas, partitioning, backup and disaster recovery, chaos, rolling deploys; abuse controls, secrets, supply-chain gates; eighteen CI jobs and two scheduled drills.}
\qa{Why it exists}{A correct system that cannot survive a lost disk, a bad deploy or a hostile client is not infrastructure.}
\qa{What it trusts}{The drills' ephemeral hosts as a lower bound on production behaviour; the operator's placement of hosts and the custody choice.}
\qa{What it refuses}{A claim without a drill; a number without a stamp; a root-running container; an unpinned image; a fixable critical CVE.}
\qa{What it can see}{Metrics, traces and alerts about the system. The correlation id on every request never touches the audit of record.}
\qa{What it cannot see}{A person: quotas, alerts and metrics are keyed by agency, route and worker.}
\qa{If it fails}{A drill that misses its ceiling fails the build; a chaos or DR finding becomes a check; a production incident is what the readiness ledger's nine operator decisions exist to prepare for.}
\qa{Checked by}{\code{check\_ha\_automation}, \code{check\_read\_replica\_routing}, \code{check\_event\_table\_partitioning}, \code{check\_dr\_drill\_scheduled}, \code{check\_chaos\_program}, \code{check\_zero\_downtime\_deploy}, \code{check\_abuse\_controls}, \code{check\_container\_hardening}, \code{check\_cve\_scanning}.}
\qa{Now or later}{Now, on drills. Multi-region, a hardware module, and production hardware are later.}
\qa{Inside and outside}{Inside: everything above. Outside: the system's view of itself, which an operator needs without a view of people.}
\end{layercard}

\nextlayer{An operator running all of this needs to see what the system is doing and to understand
where authority comes from, and must be able to do both without the system becoming a map of its
people. The next layer is how Polaris observes and understands itself.}
